\section{DESIGN OF THE INSIDE OUT FORMAT}
\label{sec:design}

\subsection{From Theory to Design}

This chapter summarizes and theoretically justifies the Inside Out format, the output of the generation and critical scrutiny phases of the RtD process. The entire first-iteration ruleset is presented in full in Annex A, and annexes B--D present its chronological evolution. The following sections argue why each structural element exists and how it serves the theoretical objectives, both for the initial design decisions made prior to playtesting and for the two significant structural additions (not just refinements of existing elements) made in light of playtest.

\subsection{Design Summary}

Inside Out is a role-playing debate format for nine participants: eight debaters organised into four teams of two, and one Evaluator. Three teams, the \textbf{Value Defenders}, are each assigned one of three values held by a fictional named Subject. The fourth team, the \textbf{Subject team}, plays that person: a specific individual facing a genuine personal dilemma, whose role is not to win but to weigh the arguments they hear and reach a decision.

Motions within this format describe a particular individual's relevant personal context, a specific situational dillema and a set of 3 values that the respective individual believes about themselves. The Value Defenders are each assigned a value and the Subject team needs to try to pick a solution for the individual described by the motion.

A debate round follows a structured sequence. The round starts with a set of Rolling In rituals. Then the three Value Defender teams deliver first speeches (\textbf{A1, B1, C1}) in turn, each arguing for their assigned value and proposing a concrete, self-constructed decision it recommends. The Subject team then delivers their own first speech (\textbf{S1}), committing to a preliminary decision and explaining which arguments have so far moved them and why. The Defenders then deliver second speeches (\textbf{A2, B2, C2}), responding to the Subject's preliminary position and refining their cases. Finally, the Subject delivers a closing synthesis (\textbf{S2}): they name the decision they have reached and explain, transparently, what arguments and values actually moved them.

A separate person, the \textbf{Evaluator}, does not judge or decide the winner. They observe the round, then award \textbf{commendations} --- named recognitions for specific demonstrated skills --- and deliver oral feedback oriented toward growth rather than competitive outcome. The round closes with \textbf{Rolling Out}: participants de-role, shake hands, and the Evaluator invites a brief, optional check on whether anything in the debate felt personally close to home.

Inside Out departs from the World Schools (WS) format in three structural ways: the \textbf{team structure and roles} (four teams with distinct functions, plus a separate Evaluator); the \textbf{shape of motions} (personal, circumstantial dilemmas rather than abstract propositions); and the \textbf{rituals and special speeches} (Rolling In/Out, the interstitial S1, and S2) that frame the debate as genuine deliberation rather than adversarial performance. The following sections justify each departure.

\subsection{Design Rationale}

Each structural choice in Inside Out traces to the theoretical framework's diagnosis of how traditional formats produce corrupted bleed, misaligned skill transfer, and container compromise. This section justifies the format's core design decisions.

\subsubsection{Why Multiple Teams Defending Values?}

Traditional formats position two teams in binary opposition, a structure whose procedural rhetoric \parencite{bogost2007persuasive} teaches that every issue has exactly two sides. Inside Out replaces this with three Value Defender teams, each advocating for a different value held by the Subject. This restructuring addresses several identified problems simultaneously.

The multi-team structure produces a procedural rhetoric of complexity and humility: the format asserts through its rules that decisions involve balancing multiple legitimate human values, not choosing between right and wrong. This directly counters the corrupted memetic bleed driving value convergence by positioning value pluralism at the core of the conflict rather than treating it as a side effect of switch-side argumentation. Because what is assigned is a value of a person (and not a direct policy solution), the design strengthens alibi \parencite{deterding2017alibis} for assuming a fragmented identity \parencite{bowman2010functions}, the debater-character relationship is explicitly about inhabiting a value, making perspective-taking the core mechanic. And because the Subject is a person with specific background, circumstances, and values rather than an abstract decision-maker, the motion reinforces that decisions are products of context, resisting the abstraction that enables generic argument templates and the drift toward what \textcite{goodnight1982spheres} identified as the technical sphere of argumentation.

Furthermore, the structure cultivates collaborative alongside combative skills. Value Defenders must persuade the Subject, a sympathetic listener making a genuine decision, not defeat an opponent. This corrects the misaligned procedural bleed by developing communication habits adapted to contexts beyond competitive debate, where persuasion typically involves collaboration rather than combat \parencite{cirlin1997sociological}. The Subject also creates a natural accessibility incentive: if a Value Defender speaks incomprehensibly, their arguments are simply not weighed in the decision, a consequence that naturally incentivises clarity.

\subsubsection{Why Split Subject and Evaluator?}

In traditional formats, a single judge both determines the winner and provides feedback, entangling evaluation with adjudication. Inside Out separates these functions: the Subject makes decisions within the debate, while the Evaluator observes and offers commendations.

This split means winning the debate is not a goal in itself but a method for growing, symbolized through commendations rather than victories. Decision-making within the debate becomes an act not of arbitrating but of synthesizing toward a position by acknowledging the interplay between values, a form of procedural rhetoric that models how real-world deliberation works, steering corrupted memetic bleed away from value convergence and toward epistemic humility. The Subject learns to synthesize; the Value Defenders learn to persuade collaboratively; and the Evaluator, freed from adjudicating competition, can focus on recognizing communicative virtue, maintaining the transformational container by decoupling status from competitive victory.

The format's final speech (S2) extends this logic to the round's conclusion. The Subject's closing synthesis requires them to name a concrete decision and explain what arguments and values actually moved them. The round ends not with one side claiming victory over the other but with an acknowledgment of how competing values were weighed and reconciled. This procedural rhetoric of synthesis over combat reinforces the format's core departure from debate-as-adversarial-validation, and models the kind of deliberative closure that characterises genuine decision-making rather than competitive performance.

\subsubsection{Why Commendations and Oral Feedback?}

Replacing scores and victory with commendations moves the format from a paradigm of competing against others toward one of improving oneself. Two debaters can be excellent in entirely different ways, a message reinforced by the diversity of commendation categories (``Compassionate Discourse'' alongside ``Strategic Rebuttal'') and the badge progression system that builds unique skill profiles over time. This directly addresses the container compromise identified in the theoretical framework: when status derives from demonstrated growth rather than competitive victory, the zero-sum dynamics that erode the transformational container \parencite{bowman2021magic, bowman2025transformative} are structurally dismantled.

The commendation system also resolves a persistent tension in modern debate formats between content density and persuasive discourse, precisely the misaligned procedural bleed identified in the theoretical framework. Commendations dissolve this conflict: debaters can pursue recognition for persuasive discourse alongside and equally with recognition for academic rigor and complexity, because these are separate categories rather than competing dimensions collapsed onto a single scale.

This growth orientation extends to post-round interaction. In traditional formats, the oral critique centres on the judge's decision: the judge explains who won and why, and advises debaters on what would have changed the result. The implicit question is ``how could I have won?'' Inside Out reframes this as oral feedback, centering on demonstrated skills and developmental trajectory. The Evaluator explains the commendations awarded, articulates why particular moments demonstrated particular skills, and offers advice oriented toward growth. The implicit question shifts to ``what am I developing, and how can I develop further?'' This connects the debate experience to the debater's broader growth trajectory rather than to a single round's competitive outcome, maintaining the transformational container by structuring the debrief around growth rather than victory.

\subsection{Design Additions from Playtesting}

Two structural additions made following the first playtest merit full theoretical justification. Both emerged from participant feedback and observational data (reported in Annex F), and both do more than refine an existing mechanism --- they introduce genuinely new procedural rhetoric. Other changes from that iteration (Points of Information, silent thinking time, motion enrichment, ritual simplification) are refinements or clarifications of mechanisms already justified above. The two additions below are not.

\subsubsection{Why the Opening Subject Speech?}

The first playtest revealed that debates began without the dilemma having been established as genuinely hard. Defenders argued into an open space and participants requested that the Subject be given a speech before Defenders begin, to give the dilemma ``real gravitas.'' The addition of an opening Subject speech (S0, here referred to as S1 in the final speech order) responds to this observation --- but its theoretical significance goes further.

By requiring the Subject to introduce who they are, what decision they face, and why no value can simply be set aside, the format ensures that the subsequent debate proceeds from a particular person's situation rather than a universal proposition. This is a structural resistance to abstraction: the question is no longer what \textit{anyone} should do about a general issue, but what \textit{this person}, with \textit{these circumstances and values}, should do in \textit{this situation}. Universalising another person's dilemma --- presuming one value or decision fits all --- becomes structurally awkward, countering what \textcite{goodnight1982spheres} identified as the drift from personal and social spheres into the technical, and what in practice manifests as an arrogant and often intolerant abstraction of singular decisions into universal prescriptions. This addition also replaces the first design's interstitial S1 (in which the Subject committed to a preliminary decision after the first Defender round) with a more theoretically motivated intervention: grounding specificity at the outset rather than mid-debate.

\subsubsection{Why the Huddle?}

The first design's interstitial S1 provided a structural check-point but did not force argumentative tensions to resolve: Defenders still argued past each other in the second round. Participants proposed adding a period in which the Subject could direct the conversation. The Huddle --- a structured interval after the first round of Defender speeches in which the Subject team may question any Defender, grant and revoke the floor, and prompt direct clashes --- replaces the old interstitial S1 entirely.

Theoretically, the Huddle simulates a council of competing interests. The dynamic is simultaneously cooperative and conflictual: Defenders can form alliances, acknowledge each other's arguments, and concede ground to gain the Subject's trust, while still competing for influence over the final decision. This captures something absent from binary debate formats: the collaborative-conflictual texture of real deliberative bodies, where stakeholders negotiate among themselves while advocating to a decision-maker. The procedural rhetoric of the Huddle teaches that resolution does not arise from one side defeating the other, but from interests being forced to confront, clarify, and --- sometimes --- reconcile.

\subsection{Synthesis and Design Tensions}

%Taken together, the design decisions map onto the three identified problems as follows. Corrupted memetic bleed (value convergence) is addressed by the multi-team structure, which positions value pluralism at the core of the conflict; by the personal Subject, which grounds debate in context-dependent reasoning; and by commendations, which sever the evaluative feedback loop. Misaligned procedural bleed (information density) is addressed by the Subject as accessibility anchor, by the collaborative persuasion dynamic, and by commendation categories that reward persuasive discourse alongside analytical rigor. Container compromise (competitive pressure) is addressed by commendations replacing victory, by the Evaluator's decoupled role, and by oral feedback oriented toward growth.

Several tensions in this design constitute productive questions for playtesting. Will competitive motivation survive without win-loss outcomes, or will debaters socialized in traditional formats experience commendations as insufficient stakes? The Subject team bears a unique cognitive burden in synthesizing three value streams, and whether this is productive challenge or overwhelming difficulty depends on experience level and motion complexity. The format's novelty (four teams, new roles, unfamiliar mechanics) may create barriers to entry for newcomers. Without victory as evaluative anchor, commendation judgments may feel more subjective; clear criteria and Evaluator training will likely be essential.